\subsection{Application to decompositions of rings}

LEMMA 2. Let A be a ring and m an ideal of A such that the ordered pair (A, m) satisfies Hensel's conditions. Let B be the quotient ring A/m and \(\pi: A \to B\) be the canonical homomorphism. For every idempotent c of B there exists a unique idempotent e of A such that \(\pi(e) = c\) .

Let \(a\) be such that \(\pi(a) = c\) ; Corollary 1 to Theorem 2 of no. 5 may be applied to the polynomial \(f(\mathbf{X}) = \mathbf{X}^2 - \mathbf{X}\) in \(\mathbf{A}[\mathbf{X}]\) and the element \(a \in \mathbf{A}\) . Then \(f'(a) = 2a - 1\) and, as \(\pi(2a - 1) = 2c - 1\) and \((2c - 1)^2 = 1\) in \(\mathbf{B}\) , \(2c - 1\) is invertible in \(\mathbf{B}\) and hence \(2a - 1\) is invertible in \(\mathbf{A}\) (§ 2, no. 13, Lemma 3). As \(f(a) \in \mathfrak{m}\) , Corollary 1 to Theorem 2 of no. 5 immediately gives the existence and uniqueness of \(e\) .

PROPOSITION 8. Let A be a ring and m an ideal of A such that the ordered pair (A, m) satisfies Hensel's conditions. Let B be the quotient ring A/m and π: A → B the canonical homomorphism. If B is the direct composition of a finite family \((\mathfrak{b}_{i})_{i\in I}\) of ideals, there exists a unique family \((\mathfrak{a}_{i})_{i\in I}\) of ideals of A such that \(\pi(\mathfrak{a}_{i}) = \mathfrak{b}_{i}\) for all \(i \in I\) and A is the direct composition of the family \((\mathfrak{a}_{i})\) .

Let \(1 = \sum_{i} c_{i}\) where \(c_{i} \in \mathfrak{b}_{i}\) for all \(i\) ; the \(c_{i}\) are idempotents of \(\mathbf{B}\) such that \(c_{i}c_{j} = 0\) for \(i \neq j\) . By Lemma 2 there therefore exist idempotents \(e_{i}\) of \(\mathbf{A} (i \in \mathbf{I})\) such that \(\pi(e_{i}) = c_{i}\) for all \(i\) ; as \(e_{i}e_{j}\) is an idempotent such that

\[
\pi (e _ {i} e _ {j}) = c _ {i} c _ {j} = 0
\]

for \(i \neq j\) , \(e_i e_j = 0\) for \(i \neq j\) (Lemma 2); as \(1 - \sum_{i} e_i\) is an idempotent such that \(\pi\left(1 - \sum_{i} e_i\right) = 1 - \sum_{i} c_i = 0\) , similarly \(1 = \sum_{i} e_i\) . It follows that \(A\) is the direct composition of the ideals \(a_i = e_i A\) and that \(\pi(a_i) = \pi(e_i) B = b_i\) .

It remains to show the uniqueness of such a decomposition. Now, suppose that A is the direct composition of another family \((\mathfrak{a}_i')_{i\in I}\) of ideals such that \(\pi (\mathfrak{a}_i^{\prime}) = \mathfrak{b}_i\) for all \(i\) ; then \(1 = \sum_{i}e_{i}^{\prime}\) where \(e_i^\prime \in \mathfrak{a}_i^\prime\) , whence, in B, \(1 = \sum_{i}\pi (e_i^{\prime})\) where \(\pi (e_i^{\prime})\in \mathfrak{b}_i\) , which implies \(\pi (e_i^{\prime}) = c_i\) ; as \(e_i^\prime\) and \(e_i\) are idempotents, necessarily \(e_i^\prime = e_i\) (Lemma 2), which completes the proof.

Remark. Proposition 8 again gives the structure of a semi-local ring on A which is Hausdorff and complete with the r-adic topology (r the Jacobson radical of A), which has already been obtained as a consequence of § 2, no. 13, Corollary to Proposition 19.

